\begin{enumerate}
\item On considère la fonction $f$ définie sur $\R$ par 
$f(x) = \e^x + x$.

\textbf{Affirmation A} : La fonction $f$ admet pour tableau de variations le tableau ci- dessous:

\begin{center}
{
\def\esp{\hspace*{5cm}}
\psset{nodesep=3pt,arrowsize=2pt 3}  
$\begin{array}{|c| *3{c}|}
\hline
 x & -\infty   & \esp & +\infty \\
 \hline
  &   &    & \Rnode{max}{+\infty}   \\
\text{variations de } f & &  &  \\
 &     \Rnode{min}{-\infty} & & 
\ncline{->}{min}{max}
\\
\hline
\end{array}$
}
\end{center}

\begin{list}{\textbullet}{}
\item $\ds\lim_{x\to -\infty} \e^{x}=0$ et $\ds\lim_{x\to -\infty} x = -\infty$ donc $\ds\lim_{x\to -\infty} f(x)=-\infty$
\item $\ds\lim_{x\to +\infty} \e^{x}=+\infty$ et $\ds\lim_{x\to +\infty} x = +\infty$ donc $\ds\lim_{x\to +\infty} f(x)=+\infty$
\item $f'(x)=\e^{x}+1>0$ sur $\R$, donc la fonction $f$ est strictement croissante sur $\R$.
\end{list}

\begin{flushright}
\textbf{\blue Affirmation A vraie}
\end{flushright}

\textbf{Affirmation B} : L'équation $f(x) = -2$ admet deux solutions dans $\R$.

\begin{center}
{
\def\esp{\hspace*{5cm}}
\psset{nodesep=3pt,arrowsize=2pt 3}  
$\begin{array}{|c| *3{c}|}
\hline
 x & -\infty   & \esp & +\infty \\
 \hline
  &   &    & \Rnode{max}{+\infty}   \\
\text{variations de } f & &  &  \\
 &     \Rnode{min}{-\infty} & & 
\ncline{->}{min}{max}
\rput*(-3,0.6){\Rnode{zero}{\red -2}}
\rput(-3,1.55){\Rnode{alpha}{\red \alpha}}
\ncline[linestyle=dotted, linecolor=red]{alpha}{zero}
\\
\hline
\end{array}$
}
\end{center}

La fonction $f$ est dérivable donc continue sur $\R$. Elle est strictement croissante sur $\R$ et va de $-\infty$ à $+\infty$. Donc, d'après le corollaire du théorème des valeurs intermédiaires, l'équation $f(x)=-2$ admet une solution unique dans $\R$.

\begin{flushright}
\textbf{\blue Affirmation B fausse}
\end{flushright}

\item \textbf{Affirmation C} :
$\displaystyle\lim_{x \to + \infty} \dfrac{\ln(x) - x^2 +2}{3x^2} = - \dfrac13$.

Pour $x\neq 0$, on a:
$\dfrac{\ln(x) - x^2 +2}{3x^2}
= \dfrac{x^2 \left (\frac{\ln(x)}{x^2} -1 +\frac{2}{x^2} \right )}{3x^2}
= -\dfrac{1}{3} \left ( 1 -  \dfrac{\ln(x)}{x^2} -\dfrac{2}{x^2}  \right )$ 

\begin{list}{\textbullet}{}
\item $\dfrac{\ln(x)}{x^2}=\dfrac{\ln(x)}{x}\times \dfrac{1}{x}$

On sait que $\ds\lim_{x\to +\infty} \dfrac{\ln(x)}{x}=0$ et que $\ds\lim_{x\to +\infty} \dfrac{1}{x}=0$; donc par produit: $\ds\lim_{x\to +\infty}\dfrac{\ln(x)}{x^2}=0$.

\item On sait aussi que $\ds\lim_{x\to +\infty} \dfrac{1}{x^2}=0$ donc $\ds\lim_{x\to +\infty} \dfrac{2}{x^2}=0$.
\end{list}

On en déduit que
$\ds\lim_{x\to +\infty} \left ( 1 -  \dfrac{\ln(x)}{x^2} -\dfrac{2}{x^2}  \right )=1$ et donc que $\displaystyle\lim_{x \to + \infty} \dfrac{\ln(x) - x^2 +2}{3x^2} = - \dfrac13$.

\begin{flushright}
\textbf{\blue Affirmation C vraie}
\end{flushright}

\item On considère la fonction $k$ définie et continue sur $\R$ par
$k(x) = 1 + 2\e^{-x^2 + 1}$.

\textbf{Affirmation D} : Il existe une primitive de la fonction $k$ décroissante sur $\R$.

Toute primitive $K$ de la fonction $k$ a pour dérivée $k$. Or, pour tout réel $X$, on a $\e^{X}>0$. Donc pour tout réel $x$, on a $\e^{-x^2 + 1}>0$, donc $1+2\e^{-x^2 + 1}>0$, et donc $k(x)>0$.

La primitive $K$ a donc une dérivée toujours strictement positive, donc elle est strictement croissante sur $\R$.

\begin{flushright}
\textbf{\blue Affirmation D fausse}
\end{flushright}

\item On considère l'équation différentielle
$\quad  (E): \quad 3y' + y = 1$.

\textbf{Affirmation E} : La fonction $g$ définie sur $\R$ par 
$g(x) = 4\e^{- \frac13 x} + 1$
est solution de l'équation différentielle $(E)$ avec $g(0) = 5$.

\begin{list}{\textbullet}{}
\item $g(x) = 4\e^{- \frac13 x} + 1$ donc $g(0)=4\e^{0} + 1=4+1=5$
\item $g$ est dérivable sur $\R$ et $g'(x)= 4\times \left (-\dfrac{1}{3}\right )\e^{-\frac{1}{3}x} = -\dfrac{4}{3}\e^{-\frac{1}{3}x} $.

Donc $3g'(x)+g(x)
= 3\times \left ( -\dfrac{4}{3}\e^{-\frac{1}{3}x}\right ) +  \left ( 4\e^{-\frac{1}{3}x} +1 \right )
= - 4\e^{-\frac{1}{3}x} + 4\e^{-\frac{1}{3}x} +1 = 1$

Donc la fonction $g$ est solution de l'équation différentielle $(E)$.
\end{list}

\begin{flushright}
\textbf{\blue Affirmation E vraie}
\end{flushright}

\item \textbf{Affirmation F} : Une intégration par parties permet d'obtenir : 
$\displaystyle\int_0^1 x\e^{-x}\:\text{d}x = 1 - 2\e^{-1}$.

En prenant:
$\left\{
\begin{array}{l !{=} l}
u(x) & x \\
v'(x)  &  \e^{-x}
\end{array}\right.$, on a:
$\left\{
\begin{array}{l !{=} l}
u'(x) & 1\\
v(x) & - \e^{-x}
\end{array}\right.$

Cela donne par intégration par parties :

$\displaystyle\int_0^1 x\e^{-x}\:\text{d}x  = \left[- x \e^{-x}\right]_0^1  - \displaystyle\int_0^1 - \e^{-x}\:\text{d}x = \left[- x \e^{-x} - \e^{-x}\right]_0^1 = - \e^{-1} - \e^{-1}  + 1 = 1 - 2\e^{- 1}$. 

\begin{flushright}
\textbf{\blue Affirmation F vraie}
\end{flushright}

\end{enumerate}

\bigskip
